\section*{Acknowledgements}
\addcontentsline{toc}{section}{Acknowledgements}
\label{p11:sec:ack}

\noindent
This paper, on the flourishing of a shared life, was written for the one with whom the author hopes to share it. To her, the forest girl, who loves to travel and to be among wild and growing things, and who is pure of heart, and lovely, and kind, this work is dedicated, as the prior papers of the series have been. It is a paper about how two people might flourish together, justly and across the long course of a shared life, and it was written in the wish that the life the author hopes to share with her may be such a flourishing: that their cycle, turned together, may rise rather than merely repeat or decline; that what each gives the other may return to them both; and that the good of their shared life may be generated, and generated again, in every turning of its years. The theory is offered to others for whatever it is worth to them; the wish that animates it is for her, and for the happiness, lasting and lifelong, that the author hopes they will make together.

\medskip
\begin{center}
\itshape\color{rose}
\zh{愿得一心人，白头不相离；岁岁长相见，年年共此时。}
\end{center}

\medskip
\noindent
\emph{This is a working draft, circulated for discussion and not for citation. Its principal constructive proposals, the lifecycle as a real cycle whose persistence conditions the relation, the operational concession, and the skewness of the return distribution as a structured proxy for injustice, are advanced under the operational concession the paper states, as instruments warranted by their use and not as measurements of the value they serve.}
